% Artículo del proyecto final. Formato NeurIPS, máximo 4 páginas de contenido.
% Las referencias, la lista de verificación y los apéndices no cuentan.
% El artículo se escribe en inglés.
%
% Los comentarios en español explican qué va en cada sección. Borrar el texto
% gris (\todo{...}) y sustituirlo por el contenido del equipo.
% Compilar con: pdflatex articulo; bibtex articulo; pdflatex articulo; pdflatex articulo
\documentclass{article}
\usepackage[preprint]{neurips_2025}  % preprint: muestra autores, sin números de línea
\usepackage[utf8]{inputenc}
\usepackage[T1]{fontenc}
\usepackage{hyperref,url,booktabs,amsfonts,amsmath,nicefrac,microtype,graphicx,xcolor}
\newcommand{\todo}[1]{\textcolor{gray}{\emph{#1}}}

\title{\todo{A title that states the main finding}}

\author{%
  \todo{Author One} \\
  \todo{Program, Institution} \\
  \texttt{\todo{email}} \\
  \And
  \todo{Author Two} \\
  \todo{Program, Institution} \\
  \texttt{\todo{email}} \\
}

\begin{document}
\maketitle

% ---------------------------------------------------------------------------
% RESUMEN. Se escribe primero, con la estructura Y, PERO, POR LO TANTO.
% 1. Y: qué se sabe y por qué importa el problema.
% 2. PERO ("However,"): qué queda sin resolver, como una pregunta concreta.
% 3. POR LO TANTO ("In this work, we"): qué hace este trabajo y con qué datos.
% 4. Resultados con números, uno por oración, contra la línea base.
% 5. Qué significa el resultado.
% ---------------------------------------------------------------------------
\begin{abstract}
  \todo{Context sentence. However, gap sentence. In this work, we ... We evaluated ... on ...
  Our method achieved X on metric Y, compared with Z for the baseline. These results suggest ...}
\end{abstract}

% ---------------------------------------------------------------------------
% INTRODUCCIÓN. Expande el resumen en el mismo orden, uno o dos párrafos por
% movimiento. Termina con la lista de contribuciones y un párrafo que dice cómo
% se organiza el resto del artículo.
% ---------------------------------------------------------------------------
\section{Introduction}
\todo{Paragraph 1 (And): the problem, who has it, and what it costs. Cite evidence.}

\todo{Paragraph 2 (But): what current approaches leave open, stated as a concrete question.}

\todo{Paragraph 3 (Therefore): what this work does. Start with ``In this work, we''.}

The contributions of this work are:
\begin{itemize}
  \item \todo{Contribution 1, one sentence.}
  \item \todo{Contribution 2, one sentence.}
\end{itemize}

\todo{Section~\ref{sec:related} reviews related work, Section~\ref{sec:method} describes the method,
Section~\ref{sec:experiments} reports the experiments, and Section~\ref{sec:limitations} discusses
limitations.}

% ---------------------------------------------------------------------------
% TRABAJO RELACIONADO. Agrupar por enfoque y decir en qué se distingue este
% trabajo de cada grupo. Cada referencia se verifica (DOI o enlace) antes de
% citarla. Una referencia inventada descuenta puntos.
% ---------------------------------------------------------------------------
\section{Related work}
\label{sec:related}
\todo{Group 1 of approaches and how this work differs. Group 2 and how this work differs.}
Clear structure helps readers follow the argument~\citep{mensh2017ten}.

% ---------------------------------------------------------------------------
% MÉTODO. Suficiente detalle para reproducir el trabajo: datos, preprocesamiento,
% modelo, hiperparámetros y cómo se eligieron. Usar tiempo pasado.
% ---------------------------------------------------------------------------
\section{Method}
\label{sec:method}
\todo{Problem formulation with notation defined at first use.}

\todo{Data: source, size, license, train/validation/test split.}

\todo{Model and training procedure.}

% ---------------------------------------------------------------------------
% EXPERIMENTOS. La línea base es obligatoria. Reportar la variabilidad
% (desviación estándar sobre varias ejecuciones independientes, cada una con
% una semilla distinta, o intervalos de confianza).
% Abrir cada subsección con el resultado principal.
% ---------------------------------------------------------------------------
\section{Experiments}
\label{sec:experiments}
\todo{Setup: baseline, metric, number of runs, hardware.}

\begin{table}[t]
  \caption{\todo{Self-contained caption: what the table shows, the metric, and whether higher is better.
  Mean and standard deviation over five independent executions.}}
  \label{tab:results}
  \centering
  \begin{tabular}{lcc}
    \toprule
    Method & \todo{Metric} $\uparrow$ & \todo{Time (s)} $\downarrow$ \\
    \midrule
    \todo{Baseline} & \todo{0.00 $\pm$ 0.00} & \todo{0.0} \\
    \todo{Ours}     & \todo{0.00 $\pm$ 0.00} & \todo{0.0} \\
    \bottomrule
  \end{tabular}
\end{table}

\todo{Main result first. Table~\ref{tab:results} shows ...}

% ---------------------------------------------------------------------------
% LIMITACIONES E IMPACTO. Qué no funciona, en qué datos podría fallar, sesgos y
% riesgos de uso. Un párrafo honesto vale más que una lista de promesas.
% ---------------------------------------------------------------------------
\section{Limitations and broader impact}
\label{sec:limitations}
\todo{Limitations of the data, the method, and the evaluation. Possible harms and mitigations.}

\section{Conclusion}
\todo{Principal finding, its implication, and the next step.}

\bibliographystyle{plainnat}
\bibliography{referencias}

% ---------------------------------------------------------------------------
% LISTA DE VERIFICACIÓN (versión breve de la NeurIPS Paper Checklist).
% Responder Yes, No o NA y justificar en una oración. No cuenta en las 4 páginas.
% ---------------------------------------------------------------------------
\appendix
\section*{Paper checklist}
\begin{enumerate}
  \item Do the abstract and introduction state the claims, and do the experiments support them?
        \todo{Yes/No/NA. Justification.}
  \item Does the paper compare against a baseline on the same data and metric?
        \todo{Yes/No/NA. Justification.}
  \item Does the paper report variability (standard deviation, confidence intervals)?
        \todo{Yes/No/NA. Justification.}
  \item Can the results be reproduced with the code and instructions in \texttt{05\_mvp/}?
        \todo{Yes/No/NA. Justification.}
  \item Are the data source, license, and preprocessing described?
        \todo{Yes/No/NA. Justification.}
  \item Does the paper discuss limitations and possible harms?
        \todo{Yes/No/NA. Justification.}
  \item Was every reference checked against its original source?
        \todo{Yes/No/NA. Justification.}
  \item Is the use of generative AI declared in \texttt{06\_ia/}?
        \todo{Yes/No/NA. Justification.}
\end{enumerate}

\end{document}
